\SourceLeaf{26}
10. Ist $\xi$ ein primes Polynom mod. $(1+x+xx)$ und seine Determinante $=D$,
so gilt allgemein

\[
\zeta^D \equiv \zeta \quad(\mathrm{mod.}\ 1+x+xx,\ \xi).
\]

11. Quadratische, kubische Reste modulo $\ldots$

12. Ist $P$ ein kubischer Rest von $\xi$, so gilt:

\[
P^{\frac{1}{3}(D-1)} \equiv 1 \quad(\mathrm{mod.}\ 1+x+xx,\ \xi).
\]



\Needspace{8\baselineskip}\begingroup\small
\begin{center}
\textbf{BEMERKUNGEN ÜBER DIE ENTWICKLUNGEN »ZUR THEORIE DER KUBISCHEN} 
\textbf{RESTE«.}
\end{center}
\noindent\hspace*{\parindent}\textsc{Gauss} hat hier versucht, seinen dritten Beweis des
quadratischen Reziprozitätsgesetzes auf die kubischen Reste zu übertragen.
Der erste Ansatz bezieht sich allein auf den Bereich der rationalen ganzen
Zahlen und verwendet folgenden Grundgedanken:

Erfüllt die Primzahl die Bedingung $p \equiv 1 \quad(\mathrm{mod.}\ 3)$, so hat die Kongruenz $x^3 \equiv 1
\quad(\mathrm{mod.}\ p)$
drei verschiedene Wurzeln, von denen eine von $1$ verschiedene $f$ heiße.
Im System aller $p-1$ modulo $p$ inkongruenten und zu
$p$ teilerfremden ganzen Zahlen bilden dann die drei Zahlen $1, f, f^2$
bei Multiplikation eine Untergruppe vom Index $\dfrac{p-1}{3}$. Für diese denke man
sich auf eine der vielen möglichen Arten ein Repräsentantensystem $S$
gewählt, das aus den Zahlen $\alpha_1, \alpha_2,
\ldots, \alpha_{\frac{p-1}{3}}$ besteht. Kennzeichnend für diese Zahlen
ist, dass die Kongruenz $\alpha_i \equiv f^{\pm1}\alpha_k
\quad(\mathrm{mod.}\ p)$ für keine zwei Indizes $i, k$ bestehen kann.
Definiert man nun die Systeme $S'$ und $S''$ mit je $\dfrac{p-1}{3}$ Zahlen
$\beta, \gamma$ durch

\[
\beta_i \equiv f\alpha_i, \qquad \gamma_i \equiv f^2\alpha_i \quad(\mathrm{mod.}\ p),
\]
\noindent so bilden die $(p-1)$ Zahlen $\alpha, \beta, \gamma$ genau das vollständige
Restsystem\hfill\\
$1, 2, \ldots, {p-1} \quad(\mathrm{mod.}\ p)$.

Ist $a$ eine beliebige durch $p$ nicht teilbare Zahl, so ist

\[
a^{\frac{p-1}{3}} \equiv f^{\lambda} \equiv \left[\frac{a}{p}\right] \quad(\mathrm{mod.}\ p)
\]
\noindent der »kubische Charakter« von $a$ in Bezug auf $p$, und dabei gilt
$\lambda = 0, 1$ oder $2$.

Um einen Satz über diesen kubischen Charakter zu erhalten, bilde man die
$\dfrac{p-1}{3}$ Produkte $a\alpha_1, a\alpha_2, \ldots
a\alpha_{\frac{p-1}{3}}$. Diese Zahlen bilden offenbar wieder ein
Repräsentantensystem, da aus der Kongruenz $a\alpha_i \equiv
f^{\pm1}a\alpha_k$ sofort $\alpha_i \equiv f^{\pm1}\alpha_k \quad(\mathrm{mod.}\ p)$ folgen würde.
Gehören nun $\mu$ der Zahlen $a\alpha_i$ zum System $S'$ und $\nu$
zum System $S''$, so ist ihr Produkt mod. $p$ offenbar zu
$f^{\mu+2\nu}. \alpha_1\alpha_2\ldots\alpha_{\frac{p-1}{3}}$ kongruent. Man erhält also:

\[
\left[\frac{a}{p}\right] \equiv a^{\frac{p-1}{3}} \equiv f^{\mu+2\nu} \quad(\mathrm{mod.}\ p),
\]
\noindent Das ist das genaue Analogon zu dem bekannten Lemma, das \textsc{Gauss}
seinem dritten Beweis des quadratischen Reziprozitätsgesetzes zugrunde legte.

